From - Mon May 22 18:44:36 2000
Received: from grumpy.usu.edu (grumpy.usu.edu [129.123.1.86])
	by math.usu.edu (8.9.3+Sun/8.9.1) with ESMTP id LAA08750
	for <symanzik@sunfs.math.usu.edu>; Sat, 20 May 2000 11:39:30 -0600 (MDT)
Received: from webmail.usu.edu ("port 1049"@webster.usu.edu [129.123.1.90])
 by cc.usu.edu (PMDF V5.2-32 #39375)
 with ESMTP id <01JPME6UJ6F08WW2T7@cc.usu.edu> for symanzik@sunfs.math.usu.edu;
 Sat, 20 May 2000 11:41:16 MDT
Date: Sat, 20 May 2000 11:41:16 -0600
From: weipingdeng <weipingdeng@cc.usu.edu>
Subject: Homework6 (Part 2)
Sender: weipingdeng <weipingdeng@cc.usu.edu>
To: symanzik@sunfs.math.usu.edu
Message-id: <39266481@webmail.usu.edu>
MIME-version: 1.0
X-Mailer: WebMail (Hydra) SMTP v3.61
Content-type: text/plain; charset="US-ASCII"
Content-transfer-encoding: 7bit
X-WebMail-UserID: weipingdeng
X-EXP32-SerialNo: 00002751
Content-Length: 5420
Status: RO
X-Mozilla-Status: 8001
X-Mozilla-Status2: 00000000
X-UIDL: 38baa60a000002f9

Evaluations for the 4 lectures

1. The topic is Basic Graphs. Histogram and Boxplot were talked.

First, Bruce and Glen used the web textbook: 
www.cyberk.com/1.0/a-5/index.html, and chose the Basic 3 to show us what the 
histogram is. The left picture is the known stem-leaf plot, and the right 
picture is the corresponding histogram. These two pictures show that they are 
similar. And students can add leaves to see how the corresponding bars in the 
histogram change. This picture helps students to compare the histogram to the 
known stem-leaf plot. For the example in the Uses 1, the picture shows how the 
histogram is constructed according to the table. I like this picture, because 
it makes everything clear on the graph by the blue connecting line and 
different colors. Looking at this graph, students will understand what the 
bars in the histogram stand for. It is interesting to watch changing the 
number of bins result in changing the histogram. In order to talk about the 
Boxplot, the uses 3 was shown. The stem-leaf plot, histogram, and boxplot are 
all given at the same time for same data. students can compare boxplot with 
other two plots and find the relationships, and add some data to get 
corresponding plots. I think this picture is very helpful to understand 
boxplot. Because the pictures in this web page are so helpful and impressive, 
it is definitely more effective than the pure 'chalk and textbook'.

The handouts of examples for histogram and boxplots, and some additional web 
sites will help students to understand the concept further.

Some Questions can make students think and answer actively, but it seems no 
questions asked to students in class. If students did some exercises in class 
using the web-tool, the lecture would be better.


2. The topic is summary statistics. Median and mean are talked.

In the beginning, Jingjing introduced summary statistics, including their 
applications and importance. That helped to motivate students. Qian opened the 
electronic textbook at http://surfsat.edu.au/surfstat/main/surfstat.html. At 
this page the definitions of Median and Mean are given, and some examples are 
shown how to compute them. After the concepts of the Median and Mean were 
talked respectively, the difference between them were indicated on the page by 
examples. Then the conlusion were drawn. This page is very good, because all 
the processes are available on the page, and then it's convenient for 
instructor to teach and for students to learn.

I like the exercises. In the first exercise, the applet can let students enter 
and change any values, and see what the median and mean are. The interesting 
thing is the median and mean are indicated by the red and green diamonds, 
which will move if the data are changed. Some other exercises are also given 
to students to do in class in other two web sites. That helped to check if 
students understand the concepts well. The wierd thing is that one answer is 
said  'wrong' by the page, but it is correct actually. So, instructors should 
be careful about this kind of things. It's good that Jingjing point out which 
one is the right answer immediately.


3. Confidence interval and hypothesis testing.

At web page: http://usatoday.com, some numerical statistical reports are 
available. That shows how the confidence interval is applied in our real life. 
Of course this can motivate students to study this concept. Then the 
electronic textbook hyperstat was used at http://davidmlane.com/hyperstat. In 
the unit 8, overview of the related concepts about confidence interval are in 
detail. That can help students to understand the concepts. In the next 
section, the formula and example are given about how to compute the confidence 
interval of mean, if standard diviation is known. To demonstrate this, the one 
linked Demos is used. Students can change some values of the parameters to get 
different confidence intervals, and then find some relationships. The same 
process is used to talk about Hypothesis testing in the unit 9. The detailed 
text and linked demos are very effective to help students to understand the 
concepts and computations.

In this web page, some other linked Demos have pictures which explain the 
concepts very well. But they were not shown in class. If some exercises had 
been practiced in class, that would be more effective.


4. Probability.

Roxanne talked about the basic concepts of the Probability first. She wrote 
the formulae on the board, and gave a example. The board work made everything 
very clear. Then, Ann demonstrated the concepts and rules by electronic tools 
at http://student.stat.wvu.edu/srs/modules/probedf/urn.html, and 
http://cyberk.com/1.0/b-2/exercises.html. The example in the first web page is 
interesting, because it looks like playing a game and counting the results. I 
believe students must like it. The M&M simulation in the second web page has 
colorful picture and show the probabilities. Some questions in the web based 
exercises are asked to students to compute the simple, intersection and union 
probabilities by looking at the pictures. That make students think and answer 
actively in class, and help teacher to get the feedback from students 
immediately.

The chalk and electronic tools supported each other very well in this lecture. 
If the lecture had not been talked in the rush, it would be more successful.

